\honest{Your Rationality is My Business}{Eliezer Yudkowsky}{2007}

Some readers of ``Lotteries: A Waste of Hope'' chided me for criticizing other people's decisions. The general question is: what business is it of mine if someone chooses to believe what is pleasant rather than what is true? \nb{One comment asked ``what is the standard by which you can judge someone who has the reverse preferences as behaving stupidly''.}

A snappy comeback would be: ``Why do you care whether I care whether someone else cares about the truth?'' But that is no answer. Mine is this. ``I believe that it is right and proper for me, as a human being, to have an interest in the future,'' and the human pursuit of truth, which has grown stronger over the generations, is part of that future. We are all players on that board whether we accept it or not. ``And that makes your rationality my business.'' \nb{In ``Why Truth?'' (November 2006) Yudkowsky had written that a sense of duty about truth makes one ``a lot more likely to believe that someone else should look behind the curtain too,'' and had added: ``I tend to be suspicious of morality as a motivation for rationality.'' The trouble named there was in one's own thinking; the safeguard offered here concerns how others are treated.}

Is this dangerous? Yes. People have been burned for not thinking as a priest wished. My proposal: ``Let's argue against bad ideas but not set their bearers on fire.'' \nb{Mill drew the same line in \textsc{On Liberty} (1859): such reasons are ``good reasons for remonstrating with him, or reasoning with him, or persuading him, or entreating him, but not for compelling him.''}

Some people avoid the syllogism ``I think Susie said a bad thing, therefore, Susie should be set on fire'' by holding that no one should judge anyone, ever. I name none of them. \nb{The comments that prompted the post question the author's standard or object to one judgment; none that could be found says that no one should ever judge.} I deny the ``therefore'' instead. I will argue against what Susie said, but I will not set her on fire or stop her talking by violence or regulation. Science's strange idea is that disagreements about fact are settled by experiment and mathematics, not by violence and decrees. I would extend that to a fair fight over the whole future, won by convincing people, not by burning them.

People who advocate relativism or selfishness do not appear to me to be truly relativistic or selfish. A true relativist would not judge. A truly selfish person would get on with making money instead of arguing. \nb{The main relativist positions in philosophy hold that moral truth is relative, not that no one may judge. That an egoist who argues for egoism is behaving oddly is a known objection to ethical egoism; the \textsc{Stanford Encyclopedia of Philosophy} says such worries ``do not show inconsistency.''} Rather, they have joined a side, Relativism, which aims to stop all players from making certain judgments, or Selfishness, which aims to make all players selfish. If there are any true relativists or selfish people, ``we do not hear them.'' That follows from how I have just defined them.

Last, I cannot help caring how you think, because each time someone turns away from the truth, the human story grows a little darker. ``In many cases, it is a small darkness only,'' and lying to yourself in private does less harm than public lies or burnings. But part of me mourns. So long as I only argue with your ideas, I believe it is right for me to care.
